\documentclass[11pt]{article}
\usepackage[margin=1in]{geometry}
\usepackage{amsmath,amssymb,amsthm}
\usepackage{booktabs}
\usepackage{hyperref}

\title{Technical Report: BHM Hyperprior Derivation from Prior Selection}
\author{Generated for Bezier Utility Estimation Project}
\date{January 2026}

\newcommand{\E}{\mathbb{E}}
\newcommand{\Var}{\text{Var}}

\begin{document}
\maketitle

\section{Overview}

This document describes how the Bayesian Hierarchical Model (BHM) hyperprior selection follows the methodology outlined in Section 3.2 of the paper (``Hyper-Prior Specification for the Bayesian Hierarchical Model'').

\subsection{Key Simplification}

\textbf{The central insight of this approach is a dramatic reduction in researcher degrees of freedom.}

Traditionally, specifying BHM hyperpriors for the B\'{e}zier utility function requires choosing 16 hyperprior parameters:
\begin{itemize}
  \item 4 shape parameters $(a, b, c, d)$
  \item Each with 4 hyperprior parameters: $(\mu_0, \tau_0^2, \alpha_0, \beta_0)$
  \item Total: $4 \times 4 = 16$ parameters
\end{itemize}

Under the constraint that the \textbf{unconditional BHM prior must equal the prior for a single random individual}, this reduces to at most \textbf{$2p$ decisions} (variance partition ratios $\rho_j$ and concentration parameters $\alpha_{0,j}$ for each of $p$ parameters). In practice, using common values $\rho_j = \rho$ and $\alpha_{0,j} = \alpha_0$ further reduces this to \textbf{just 2 decisions}.

\subsection{The Prior Selection Question}

The prior selection exercise answers the question: \textit{``What might the preferences of a single random individual, drawn from some population, look like?''}

This question is fundamentally about an \textbf{individual}, not about a ``representative agent'' or ``pooled'' model. However, the answer implicitly integrates over the heterogeneity in the population from which that individual was sampled:
\begin{equation}
  p(\theta) = \int_\alpha p(\theta | \alpha) p(\alpha) \, d\alpha
  \label{eq:unconditional}
\end{equation}
where $\theta$ represents the individual's parameters and $\alpha$ represents the (unknown) population distribution.

When we ask ``what do credible preferences look like for a random person?'' we are simultaneously expressing:
\begin{enumerate}
  \item Uncertainty about this particular individual's preferences
  \item Uncertainty about the population distribution they were drawn from
\end{enumerate}

The prior selection exercise captures both sources of uncertainty in a single distribution $p(\theta)$. The BHM then \textit{decomposes} this into explicit population-level and individual-level components, but the unconditional prior must remain consistent with our original answer about a single random individual.

\section{Paper Specification}

\subsection{Individual-Level Model}

The paper specifies that individual parameters follow logit-normal distributions:
\begin{align}
  a_i &\sim \text{Logit-Normal}(a_\mu, a_\sigma^2) \\
  b_i &\sim \text{Logit-Normal}(b_\mu, b_\sigma^2) \\
  c_i &\sim \text{Logit-Normal}(c_\mu, c_\sigma^2) \\
  d_i &\sim \text{Logit-Normal}(d_\mu, d_\sigma^2) \\
  \lambda_i &\sim \text{Log-Normal}(\lambda_\mu, \lambda_\sigma^2)
\end{align}
where the logit-normal distribution means:
\[
  \text{logit}(\theta_i) = \log\left(\frac{\theta_i}{1-\theta_i}\right) \sim \text{Normal}(\mu, \sigma^2)
\]

\subsection{Hyperprior Structure}

The paper specifies conjugate hyperpriors:
\begin{align}
  \theta_\mu &\sim \text{Normal}(\mu_0, \tau_0^2) \\
  \theta_{\sigma^2} &\sim \text{Inverse-Gamma}(\alpha_0, \beta_0)
\end{align}
This conjugate structure enables efficient Gibbs sampling:
\begin{itemize}
  \item Normal-Normal conjugacy for updating population means
  \item Normal-Inverse-Gamma conjugacy for updating population variances
\end{itemize}

\section{Deriving Independent Hyperpriors from Prior Selection}

This section provides a detailed derivation of the independent hyperprior structure, explaining both the mathematics and the reasoning behind each step.

\subsection{The Problem: 16 Parameters to Specify}

The independent hyperprior structure requires specifying 4 hyperparameters for each of the $p = 4$ shape parameters:
\begin{center}
\begin{tabular}{cl}
\toprule
Hyperparameter & Meaning \\
\midrule
$\mu_{0,j}$ & Prior mean of population mean $\mu_j$ \\
$\tau_{0,j}^2$ & Prior variance of population mean $\mu_j$ \\
$\alpha_{0,j}$ & Shape parameter of Inverse-Gamma prior on $\sigma_j^2$ \\
$\beta_{0,j}$ & Scale parameter of Inverse-Gamma prior on $\sigma_j^2$ \\
\bottomrule
\end{tabular}
\end{center}

This gives $4 \times 4 = 16$ hyperparameters that the researcher must specify. Without a systematic approach, this would require 16 separate judgments about interdependent quantities---an impractical and error-prone task.

\textbf{Our approach:} We impose the constraint that the \textit{unconditional} BHM prior must equal the prior selection distribution. This constraint absorbs 8 of the 16 degrees of freedom, leaving 8 remaining choices which we reparameterize as variance partition ratios $\rho_j$ and concentration parameters $\alpha_{0,j}$.

\textbf{Important caveat:} The reduction from 16 to 8 is in some sense a ``pseudo'' reduction. The prior selection exercise itself required researcher choices: specifying 4 means and 4 variances for the (independent) multivariate normal distribution over logit-transformed parameters. These 8 prior selection choices, combined with the 8 remaining BHM choices ($\rho_j$ and $\alpha_{0,j}$), bring the total researcher input back to 16.

\textbf{So what have we gained?} The value of this approach is not in reducing the total number of choices, but in \textit{where} those choices are made:
\begin{itemize}
  \item \textbf{Prior selection choices} (8): Made by reasoning about ``what might a random individual's preferences look like?'' This is cognitively tractable---we can visualize the prior predictive distribution of utility functions and RRA values, checking whether the implied preferences are reasonable.
  \item \textbf{BHM structure choices} (8): Made by reasoning about variance decomposition---how much of the prior uncertainty reflects ``not knowing the population mean'' versus ``expected heterogeneity across individuals.'' These are more abstract but have sensible defaults ($\rho = 0.3$, $\alpha_0 = 6$).
\end{itemize}

The key insight is that it is far easier to begin by reasoning about a prior predictive distribution for an individual---which can be simulated and inspected---than to directly specify 16 interdependent hyperparameters whose joint implications are difficult to assess.

\begin{center}
\begin{tabular}{lcc}
\toprule
Step & Constraint/Choice & Degrees of Freedom \\
\midrule
Starting point & --- & 16 \\
Mean constraint (Section 3.4) & $\mu_{0,j}$ determined by $\boldsymbol{\mu}_{\text{ps}}$ & $16 - 4 = 12$ \\
Variance constraint (Section 3.4) & One equation per parameter & $12 - 4 = 8$ \\
\midrule
Reparameterization & $\rho_j$ and $\alpha_{0,j}$ for each $j$ & 8 \\
Common $\rho$ (Section 3.9) & $\rho_j = \rho$ for all $j$ & $8 - 3 = 5$ \\
Common $\alpha_0$ (Section 3.9) & $\alpha_{0,j} = \alpha_0$ for all $j$ & $5 - 3 = 2$ \\
\bottomrule
\end{tabular}
\end{center}

Using common values for $\rho$ and $\alpha_0$ further reduces the 8 BHM structure choices to just 2, giving a final accounting of $8 + 2 = 10$ explicit researcher decisions (versus 16 opaque hyperparameters).

\subsection{The Starting Point: Prior Selection Results}

The prior selection exercise produces a multivariate normal distribution on the logit scale:
\[
  \begin{pmatrix} \text{logit}(a) \\ \text{logit}(b) \\ \text{logit}(c) \\ \text{logit}(d) \end{pmatrix}
  \sim \text{MVN}(\boldsymbol{\mu}_{\text{ps}}, \boldsymbol{\Sigma}_{\text{ps}})
\]
where ``ps'' denotes ``prior selection.'' This distribution represents our beliefs about what a \textit{single random individual's} preferences might look like, before seeing any data.

For the ATL dataset:
\begin{align*}
  \boldsymbol{\mu}_{\text{ps}} &= (-1.61, -1.07, 1.81, 1.98)^T
\end{align*}
\[
  \boldsymbol{\Sigma}_{\text{ps}} = \begin{pmatrix}
    0.29 & 0.08 & -0.10 & 0.05 \\
    0.08 & 0.56 & -0.21 & -0.06 \\
    -0.10 & -0.21 & 0.92 & 0.45 \\
    0.05 & -0.06 & 0.45 & 1.26
  \end{pmatrix}
\]

\subsection{The BHM Structure}

The Bayesian Hierarchical Model decomposes the individual-level distribution into two levels:

\textbf{Level 1 (Individual):} Each subject's parameter is drawn from a population distribution:
\[
  \text{logit}(\theta_{ij}) \sim \text{Normal}(\mu_j, \sigma_j^2)
\]
where $j \in \{a, b, c, d\}$ indexes the parameter and $i$ indexes the subject.

\textbf{Level 2 (Population):} The population parameters have hyperpriors:
\begin{align}
  \mu_j &\sim \text{Normal}(\mu_{0,j}, \tau_{0,j}^2) \\
  \sigma_j^2 &\sim \text{Inverse-Gamma}(\alpha_{0,j}, \beta_{0,j})
\end{align}

\subsection{The Key Constraint: Matching Unconditional Distribution}

The crucial insight is that if we integrate out the population parameters, the \textit{unconditional} distribution of $\theta_{ij}$ should match the prior selection distribution. This ensures consistency: our beliefs about a random individual are the same whether we think about them directly (prior selection) or as a draw from an unknown population (BHM).

\textbf{Matching the mean ($16 \to 12$ degrees of freedom):}

For the unconditional mean:
\[
  \E[\text{logit}(\theta_{ij})] = \E[\E[\text{logit}(\theta_{ij}) | \mu_j]] = \E[\mu_j] = \mu_{0,j}
\]
For this to equal $(\boldsymbol{\mu}_{\text{ps}})_j$, we must set $\mu_{0,j} = (\boldsymbol{\mu}_{\text{ps}})_j$. This completely determines all 4 values of $\mu_{0,j}$---no researcher choice is involved. \textbf{This eliminates 4 degrees of freedom}, reducing from 16 to 12.

\textbf{Matching the variance ($12 \to 8$ degrees of freedom):}

For the unconditional variance, we use the law of total variance:
\begin{align}
  \Var[\text{logit}(\theta_{ij})] &= \E[\Var[\text{logit}(\theta_{ij}) | \mu_j, \sigma_j^2]] + \Var[\E[\text{logit}(\theta_{ij}) | \mu_j, \sigma_j^2]] \nonumber \\
  &= \E[\sigma_j^2] + \Var[\mu_j] \nonumber \\
  &= \underbrace{\frac{\beta_{0,j}}{\alpha_{0,j} - 1}}_{\text{expected population variance}} + \underbrace{\tau_{0,j}^2}_{\text{uncertainty about pop.\ mean}}
  \label{eq:total_variance}
\end{align}

For this to equal the prior selection variance $(\boldsymbol{\Sigma}_{\text{ps}})_{jj}$, we need:
\begin{equation}
  \frac{\beta_{0,j}}{\alpha_{0,j} - 1} + \tau_{0,j}^2 = (\boldsymbol{\Sigma}_{\text{ps}})_{jj}
  \label{eq:variance_constraint}
\end{equation}

This is one constraint per parameter, involving three unknowns ($\tau_{0,j}^2$, $\alpha_{0,j}$, $\beta_{0,j}$) each. Four constraints on 12 unknowns \textbf{eliminates 4 more degrees of freedom}, leaving $12 - 4 = 8$.

\subsection{Reparameterizing the Remaining 8 Degrees of Freedom}

We now have 8 degrees of freedom: for each parameter $j$, the constraint~\eqref{eq:variance_constraint} is one equation in three unknowns ($\tau_{0,j}^2$, $\alpha_{0,j}$, $\beta_{0,j}$), leaving 2 free choices per parameter.

Rather than asking the researcher to directly specify values like $\tau_{0,j}^2 = 0.077$ or $\beta_{0,j} = 0.895$, we reparameterize in terms of quantities with clearer interpretations: the \textit{variance partition ratio} $\rho_j$ and the \textit{concentration parameter} $\alpha_{0,j}$.

\textbf{Constraint 1: Variance Partition Ratio}

We introduce a variance partition ratio $\rho_j \in (0, 1)$ for each parameter $j$, representing the fraction of total variance attributed to uncertainty about the population mean:
\begin{equation}
  \rho_j = \frac{\tau_{0,j}^2}{(\boldsymbol{\Sigma}_{\text{ps}})_{jj}}
  \label{eq:rho_def}
\end{equation}

This gives us:
\begin{align}
  \tau_{0,j}^2 &= \rho_j \cdot (\boldsymbol{\Sigma}_{\text{ps}})_{jj} \label{eq:tau_from_rho} \\
  \E[\sigma_j^2] = \frac{\beta_{0,j}}{\alpha_{0,j} - 1} &= (1 - \rho_j) \cdot (\boldsymbol{\Sigma}_{\text{ps}})_{jj} \label{eq:Esigma_from_rho}
\end{align}

\textbf{Critical observation:} For \textit{any} choice of $\rho_j \in (0, 1)$---whether different values for each parameter or a common value across all parameters---the unconditional variance of $\text{logit}(\theta_j)$ equals the prior selection variance $(\boldsymbol{\Sigma}_{\text{ps}})_{jj}$ by construction. This is because:
\[
  \Var[\text{logit}(\theta_j)] = \underbrace{(1 - \rho_j)(\boldsymbol{\Sigma}_{\text{ps}})_{jj}}_{\E[\sigma_j^2]} + \underbrace{\rho_j (\boldsymbol{\Sigma}_{\text{ps}})_{jj}}_{\tau_{0,j}^2} = (\boldsymbol{\Sigma}_{\text{ps}})_{jj}
\]
The two terms always sum to the target variance regardless of how the partition is chosen. Thus, the choice of $\rho_j$ values affects only how the variance is \textit{decomposed} between population heterogeneity and uncertainty about the population mean---it does not affect the unconditional distribution of individual parameters, which remains equivalent to the prior selection distribution by definition.

\textbf{Interpretation of $\rho_j$:}
\begin{itemize}
  \item $\rho_j \to 0$: Almost all prior uncertainty for parameter $j$ comes from population heterogeneity. We are confident about where the population is centered, but individuals vary greatly.
  \item $\rho_j \to 1$: Almost all prior uncertainty comes from not knowing the population mean. Once we learn $\mu_j$, individuals cluster tightly around it.
  \item $\rho_j = 0.3$: 30\% of prior width reflects uncertainty about the population center; 70\% reflects expected individual heterogeneity.
\end{itemize}

\textbf{Constraint 2: Concentration Parameter $\alpha_0$}

Equation~\eqref{eq:Esigma_from_rho} still has two unknowns: $\alpha_{0,j}$ and $\beta_{0,j}$. We need one more constraint.

The Inverse-Gamma distribution has mean and variance:
\begin{align}
  \E[\sigma^2] &= \frac{\beta_0}{\alpha_0 - 1} \quad \text{for } \alpha_0 > 1 \\
  \Var[\sigma^2] &= \frac{\beta_0^2}{(\alpha_0 - 1)^2 (\alpha_0 - 2)} \quad \text{for } \alpha_0 > 2
\end{align}

The coefficient of variation (CV) of $\sigma^2$ is:
\begin{equation}
  \text{CV}[\sigma^2] = \frac{\sqrt{\Var[\sigma^2]}}{\E[\sigma^2]} = \frac{1}{\sqrt{\alpha_0 - 2}}
  \label{eq:cv_ig}
\end{equation}

The CV measures how uncertain we are about the population variance \textit{relative to its expected value}. A small CV means we're confident the true $\sigma^2$ is close to $\E[\sigma^2]$; a large CV means $\sigma^2$ could be much smaller or larger than expected.

\textbf{Why CV = 0.5?} This is a modeling choice reflecting moderate uncertainty about the population variance. Let's consider the implications:
\begin{center}
\begin{tabular}{ccc}
\toprule
CV & $\alpha_0$ & Interpretation \\
\midrule
1.0 & 3 & High uncertainty: $\sigma^2$ could easily be 2$\times$ or 0.5$\times$ the expected value \\
0.71 & 4 & Moderate-high uncertainty \\
0.5 & 6 & Moderate uncertainty: $\pm$50\% variation is plausible \\
0.33 & 11 & Low uncertainty: $\sigma^2$ is fairly well-determined \\
0.25 & 18 & Very low uncertainty: strong prior belief about $\sigma^2$ \\
\bottomrule
\end{tabular}
\end{center}

With CV = 0.5 ($\alpha_0 = 6$):
\begin{itemize}
  \item The prior allows $\sigma^2$ to range from roughly 0.3$\times$ to 2$\times$ its expected value with reasonable probability
  \item We're not dogmatic about the prior selection variance being exactly correct
  \item But we're not so uncertain that the data has no prior guidance
\end{itemize}

\textbf{Alternative interpretation via effective sample size:} For the Inverse-Gamma, $\alpha_0$ can be interpreted as an effective sample size---the number of ``prior observations'' we have about $\sigma^2$. With $\alpha_0 = 6$, we have the equivalent of 6 prior observations, which provides moderate information while leaving substantial room for the data to inform the posterior.

\subsection{Solving for the Hyperprior Parameters}

Given the prior selection results $(\boldsymbol{\mu}_{\text{ps}}, \boldsymbol{\Sigma}_{\text{ps}})$, a variance partition ratio $\rho_j \in (0,1)$, and a concentration parameter $\alpha_{0,j} > 2$ for each parameter $j$:

\textbf{Step 1: Mean hyperprior location}
\[
  \mu_{0,j} = (\boldsymbol{\mu}_{\text{ps}})_j
\]

\textbf{Step 2: Mean hyperprior variance (from $\rho_j$)}
\[
  \tau_{0,j}^2 = \rho_j \cdot V_j \quad \text{where } V_j = (\boldsymbol{\Sigma}_{\text{ps}})_{jj}
\]

\textbf{Step 3: Expected population variance (from $\rho_j$)}
\[
  \E[\sigma_j^2] = (1 - \rho_j) \cdot V_j
\]

\textbf{Step 4: Inverse-Gamma scale parameter (from $\alpha_{0,j}$ and $\E[\sigma_j^2]$)}
\[
  \beta_{0,j} = \E[\sigma_j^2] \cdot (\alpha_{0,j} - 1) = (1 - \rho_j) \cdot V_j \cdot (\alpha_{0,j} - 1)
\]

\subsection{Verification: The Unconditional Variance Matches}

We can verify that our construction satisfies the constraint~\eqref{eq:variance_constraint} for any valid $\rho_j$:
\begin{align*}
  \frac{\beta_{0,j}}{\alpha_{0,j} - 1} + \tau_{0,j}^2 &= \frac{(1-\rho_j) V_j (\alpha_{0,j} - 1)}{\alpha_{0,j} - 1} + \rho_j V_j \\
  &= (1 - \rho_j) V_j + \rho_j V_j \\
  &= V_j = (\boldsymbol{\Sigma}_{\text{ps}})_{jj} \quad \checkmark
\end{align*}
where $V_j = (\boldsymbol{\Sigma}_{\text{ps}})_{jj}$ is the target variance from prior selection.

\textbf{This identity holds for any $\rho_j \in (0, 1)$ and any $\alpha_{0,j} > 2$.} Whether the researcher chooses different partition ratios for each parameter (e.g., $\rho_a = 0.2$, $\rho_b = 0.4$, $\rho_c = 0.3$, $\rho_d = 0.1$) or a common value ($\rho_j = 0.3$ for all $j$), the unconditional distribution of each parameter matches the prior selection target by construction. The choice of partition ratios is purely about expressing beliefs regarding the \textit{source} of variance---not its magnitude.

\subsection{What About Correlations?}

The independent structure ignores the off-diagonal elements of $\boldsymbol{\Sigma}_{\text{ps}}$. Under this structure:
\[
  \text{Cov}[\text{logit}(\theta_j), \text{logit}(\theta_k)] = 0 \quad \text{for } j \neq k
\]

This is a limitation. The prior selection found $\text{Corr}[a, c] = -0.35$, which encodes the constraint that utility functions with low initial curvature tend to have high upper curvature. The independent structure loses this information.

However, with sufficient data (120 choices per subject), the likelihood will effectively rule out implausible $(a, c)$ combinations, mitigating the impact of this approximation. For small datasets, the correlated (NIW) structure may be preferable despite its computational challenges.

\subsection{Reducing Degrees of Freedom: Common $\rho$ and $\alpha_0$ ($8 \to 2$)}

At this point, the prior selection constraint has reduced 16 hyperparameters to 8 researcher choices: $\rho_j$ and $\alpha_{0,j}$ for each of the $p = 4$ parameters. This flexibility is available because, as shown above, \textit{any} valid choice of these 8 values produces an unconditional distribution matching the prior selection target.

However, we rarely have strong prior beliefs that would justify different partition ratios or concentration parameters across the Bézier shape parameters. Absent such beliefs, we set:
\begin{align}
  \rho_j &= \rho \quad \text{for all } j \quad \text{(reduces 4 choices to 1)} \\
  \alpha_{0,j} &= \alpha_0 \quad \text{for all } j \quad \text{(reduces 4 choices to 1)}
\end{align}

This reduces the researcher degrees of freedom from 8 to 2: a single $\rho$ and a single $\alpha_0$. The complete accounting:
\[
  \underbrace{16}_{\text{starting}} \xrightarrow{\text{mean constraint}} \underbrace{12}_{\mu_{0,j} \text{ fixed}} \xrightarrow{\text{variance constraint}} \underbrace{8}_{\text{reparameterized}} \xrightarrow{\text{common } \rho} \underbrace{5}_{} \xrightarrow{\text{common } \alpha_0} \underbrace{2}_{\text{final}}
\]

\textbf{When might different $\rho_j$ values be warranted?}

A researcher might choose different partition ratios if they have prior beliefs about how variance is distributed differently across parameters. For example:
\begin{itemize}
  \item If prior literature strongly establishes the population mean of parameter $a$ (implying low $\rho_a$) but says little about $c$ (implying higher $\rho_c$)
  \item If some parameters are known to be highly heterogeneous across populations while others are relatively stable
\end{itemize}

Absent such specific prior information, using a common $\rho$ is more parsimonious and easier to interpret: the statement ``30\% of prior variance reflects uncertainty about population location'' applies uniformly across all parameters.

\textbf{Further simplification:} In most applications, we also hold $\alpha_0$ constant (e.g., $\alpha_0 = 6$ for all parameters), reducing the researcher's task to selecting a single value $\rho$. This single parameter has a clear interpretation: the fraction of total prior variance attributed to not knowing where the population is centered.

\subsection{Summary of the Independent Hyperprior Derivation}

\begin{center}
\begin{tabular}{ll}
\toprule
\textbf{Given} & \textbf{From} \\
\midrule
$\boldsymbol{\mu}_{\text{ps}}$, $\boldsymbol{\Sigma}_{\text{ps}}$ & Prior selection exercise \\
$\rho_j \in (0, 1)$ for each $j$ & Researcher choice (we use common $\rho = 0.3$) \\
$\alpha_{0,j} > 2$ for each $j$ & Researcher choice (we use common $\alpha_0 = 6$) \\
\midrule
\textbf{Derived} & \textbf{Formula} \\
\midrule
$\mu_{0,j}$ & $(\boldsymbol{\mu}_{\text{ps}})_j$ \\
$\tau_{0,j}^2$ & $\rho_j \cdot (\boldsymbol{\Sigma}_{\text{ps}})_{jj}$ \\
$\beta_{0,j}$ & $(1 - \rho_j) \cdot (\boldsymbol{\Sigma}_{\text{ps}})_{jj} \cdot (\alpha_{0,j} - 1)$ \\
\bottomrule
\end{tabular}
\end{center}

\textbf{Key property:} For any valid choice of $\rho_j$ and $\alpha_{0,j}$ values, the unconditional distribution of $\theta_j$ matches the prior selection target by construction.

\textbf{Practical simplification:} While the general formulation allows $2p$ researcher choices ($\rho_j$ and $\alpha_{0,j}$ for each of $p$ parameters), using common values $\rho_j = \rho$ and $\alpha_{0,j} = \alpha_0$ reduces this to just 2 choices. This is our approach, absent specific prior beliefs about differential variance decomposition across parameters.

\section{Derived Hyperpriors for ATL Dataset}

Applying the above methodology to the ATL prior selection results:

\begin{table}[h]
\centering
\begin{tabular}{lccccc}
\toprule
Parameter & $V_{\text{target}}$ & $\mu_0$ & $\tau_0^2$ & $\alpha_0$ & $\beta_0$ \\
\midrule
$a$ & 0.256 & $-1.641$ & 0.077 & 6.0 & 0.895 \\
$b$ & 0.096 & $-0.796$ & 0.029 & 6.0 & 0.338 \\
$c$ & 5.813 & $4.670$ & 1.744 & 6.0 & 20.34 \\
$d$ & 6.452 & $4.438$ & 1.936 & 6.0 & 22.58 \\
$\mu$ & 1.000 & $-2.300$ & 0.300 & 6.0 & 3.500 \\
\bottomrule
\end{tabular}
\caption{Derived BHM hyperpriors from ATL prior selection}
\end{table}

\subsection{Verification of Unconditional Variance}

For each parameter, we verify that the unconditional variance matches the target:
\[
  \E[\sigma^2] + \tau_0^2 = \frac{\beta_0}{\alpha_0 - 1} + \tau_0^2 = \frac{\beta_0}{5} + \tau_0^2
\]

For parameter $a$:
\[
  \frac{0.895}{5} + 0.077 = 0.179 + 0.077 = 0.256 \quad \checkmark
\]

\section{Deriving Correlated (NIW) Hyperpriors from Prior Selection}

This section provides a detailed derivation of the correlated hyperprior structure using the Normal-Inverse-Wishart (NIW) conjugate prior, paralleling the development in Section 3.

\subsection{The Problem: 16 Parameters to Specify}

The correlated hyperprior structure requires specifying parameters for the Normal-Inverse-Wishart prior:
\begin{center}
\begin{tabular}{clc}
\toprule
Hyperparameter & Meaning & Count \\
\midrule
$\boldsymbol{\mu}_0$ & Prior mean vector for population mean & 4 \\
$\kappa_0$ & Shrinkage parameter (how tightly $\mu$ is centered on $\mu_0$) & 1 \\
$\boldsymbol{\Psi}_0$ & Scale matrix for Inverse-Wishart (symmetric $4 \times 4$) & 10 \\
$\nu_0$ & Degrees of freedom for Inverse-Wishart & 1 \\
\midrule
& \textbf{Total} & \textbf{16} \\
\bottomrule
\end{tabular}
\end{center}

As with the independent case, directly specifying 16 interdependent parameters is impractical. We use the same approach: impose the constraint that the unconditional BHM prior must equal the prior selection distribution, then reparameterize the remaining degrees of freedom.

\textbf{Important:} The NIW structure has exactly the same number of hyperparameters (16) as the independent structure, but they are organized differently---a 4-vector, two scalars, and a $4 \times 4$ symmetric matrix rather than 4 groups of 4 scalars.

\begin{center}
\begin{tabular}{lcc}
\toprule
Step & Constraint/Choice & Degrees of Freedom \\
\midrule
Starting point & --- & 16 \\
Mean constraint (Section 5.4) & $\boldsymbol{\mu}_0 = \boldsymbol{\mu}_{\text{ps}}$ & $16 - 4 = 12$ \\
Covariance constraint (Section 5.4) & Matrix equation (10 unique elements) & $12 - 10 = 2$ \\
\midrule
Reparameterization & $\rho$ and $\nu_0$ & 2 \\
\bottomrule
\end{tabular}
\end{center}

\textbf{Key difference from independent case:} The covariance constraint is a \textit{matrix} equation with 10 independent constraints (for the symmetric $4 \times 4$ matrix), which absorbs 10 degrees of freedom at once. This leaves only 2 remaining choices, compared to 8 in the independent case.

\textbf{Pseudo-reduction caveat:} As before, the prior selection exercise required specifying a mean vector (4 values) and covariance matrix (10 unique values), totaling 14 choices. Combined with the 2 NIW structure choices ($\rho$ and $\nu_0$), the total researcher input is 16---the same as directly specifying hyperparameters. The value is in \textit{where} choices are made: prior selection choices are cognitively tractable (visualize utility functions), while $\rho$ and $\nu_0$ have clear interpretations.

\subsection{The Starting Point: Prior Selection Results}

The prior selection exercise produces a multivariate normal distribution on the logit scale:
\[
  \boldsymbol{\theta} = \begin{pmatrix} \text{logit}(a) \\ \text{logit}(b) \\ \text{logit}(c) \\ \text{logit}(d) \end{pmatrix}
  \sim \text{MVN}(\boldsymbol{\mu}_{\text{ps}}, \boldsymbol{\Sigma}_{\text{ps}})
\]

Unlike the independent case, we now use the \textit{full} covariance matrix $\boldsymbol{\Sigma}_{\text{ps}}$, including off-diagonal correlations. For the ATL dataset:
\[
  \boldsymbol{\Sigma}_{\text{ps}} = \begin{pmatrix}
    0.29 & 0.08 & -0.10 & 0.05 \\
    0.08 & 0.56 & -0.21 & -0.06 \\
    -0.10 & -0.21 & 0.92 & 0.45 \\
    0.05 & -0.06 & 0.45 & 1.26
  \end{pmatrix}
\]

The correlation between $a$ and $c$ is approximately $-0.19$, encoding the constraint that utility functions with certain initial curvature tend to have compensating upper curvature.

\subsection{The NIW BHM Structure}

The Normal-Inverse-Wishart conjugate prior generalizes the univariate Normal-Inverse-Gamma:

\textbf{Level 1 (Individual):} Each subject's parameters are drawn from a population distribution:
\[
  \boldsymbol{\theta}_i | \boldsymbol{\mu}, \boldsymbol{\Sigma} \sim \text{MVN}(\boldsymbol{\mu}, \boldsymbol{\Sigma})
\]

\textbf{Level 2 (Population):} The population parameters have NIW hyperpriors:
\begin{align}
  \boldsymbol{\mu} | \boldsymbol{\Sigma} &\sim \text{MVN}(\boldsymbol{\mu}_0, \boldsymbol{\Sigma}/\kappa_0) \\
  \boldsymbol{\Sigma} &\sim \text{Inverse-Wishart}(\nu_0, \boldsymbol{\Psi}_0)
\end{align}

The conditional structure $\boldsymbol{\mu} | \boldsymbol{\Sigma}$ is key: uncertainty about the population mean scales with the population covariance, preserving the correlation structure at both levels.

\subsection{The Key Constraint: Matching Unconditional Distribution}

\textbf{Matching the mean ($16 \to 12$ degrees of freedom):}

The unconditional mean is:
\[
  \E[\boldsymbol{\theta}_i] = \E[\E[\boldsymbol{\theta}_i | \boldsymbol{\mu}]] = \E[\boldsymbol{\mu}] = \boldsymbol{\mu}_0
\]
For this to equal $\boldsymbol{\mu}_{\text{ps}}$, we must set $\boldsymbol{\mu}_0 = \boldsymbol{\mu}_{\text{ps}}$. This completely determines the 4-element mean vector---no researcher choice is involved. \textbf{This eliminates 4 degrees of freedom.}

\textbf{Matching the covariance ($12 \to 2$ degrees of freedom):}

The unconditional covariance uses the law of total covariance:
\begin{align}
  \Var[\boldsymbol{\theta}_i] &= \E[\Var[\boldsymbol{\theta}_i | \boldsymbol{\mu}, \boldsymbol{\Sigma}]] + \Var[\E[\boldsymbol{\theta}_i | \boldsymbol{\mu}, \boldsymbol{\Sigma}]] \nonumber \\
  &= \E[\boldsymbol{\Sigma}] + \Var[\boldsymbol{\mu}] \nonumber \\
  &= \underbrace{\frac{\boldsymbol{\Psi}_0}{\nu_0 - p - 1}}_{\text{expected population cov.}} + \underbrace{\frac{1}{\kappa_0} \cdot \frac{\boldsymbol{\Psi}_0}{\nu_0 - p - 1}}_{\text{uncertainty about pop.\ mean}}
  \label{eq:niw_total_cov}
\end{align}

This simplifies to:
\begin{equation}
  \Var[\boldsymbol{\theta}_i] = \left(1 + \frac{1}{\kappa_0}\right) \cdot \frac{\boldsymbol{\Psi}_0}{\nu_0 - p - 1}
  \label{eq:niw_unconditional_cov}
\end{equation}

For this to equal $\boldsymbol{\Sigma}_{\text{ps}}$, we need:
\begin{equation}
  \boldsymbol{\Psi}_0 = \frac{\kappa_0}{\kappa_0 + 1} \cdot (\nu_0 - p - 1) \cdot \boldsymbol{\Sigma}_{\text{ps}}
  \label{eq:psi_constraint}
\end{equation}

This is a \textit{matrix} equation. Given choices of $\kappa_0$ and $\nu_0$, the entire $4 \times 4$ scale matrix $\boldsymbol{\Psi}_0$ is determined (10 unique elements). \textbf{This eliminates 10 degrees of freedom}, leaving only 2.

\subsection{Reparameterizing the Remaining 2 Degrees of Freedom}

The 2 remaining degrees of freedom are $\kappa_0$ and $\nu_0$. We reparameterize $\kappa_0$ in terms of the variance partition ratio $\rho$ for interpretability.

\textbf{Researcher Choice 1: Variance Partition Ratio $\rho$}

Define $\rho$ as the fraction of unconditional variance from uncertainty about the population mean:
\[
  \rho = \frac{\Var[\boldsymbol{\mu}]}{\Var[\boldsymbol{\theta}]} = \frac{\frac{1}{\kappa_0} \E[\boldsymbol{\Sigma}]}{\left(1 + \frac{1}{\kappa_0}\right) \E[\boldsymbol{\Sigma}]} = \frac{1}{1 + \kappa_0}
\]

Solving for $\kappa_0$:
\begin{equation}
  \kappa_0 = \frac{1 - \rho}{\rho}
  \label{eq:kappa_from_rho}
\end{equation}

\begin{center}
\begin{tabular}{ccc}
\toprule
$\rho$ & $\kappa_0$ & Interpretation \\
\midrule
0.1 & 9.0 & 10\% mean uncertainty, 90\% heterogeneity \\
0.3 & 2.33 & 30\% mean uncertainty, 70\% heterogeneity (our choice) \\
0.5 & 1.0 & Equal split \\
0.7 & 0.43 & 70\% mean uncertainty, 30\% heterogeneity \\
\bottomrule
\end{tabular}
\end{center}

\textbf{Researcher Choice 2: Degrees of Freedom $\nu_0$}

The degrees of freedom $\nu_0$ controls how concentrated the Inverse-Wishart distribution is around its expected value. Must satisfy $\nu_0 > p + 1 = 5$ for finite variance.

\begin{center}
\begin{tabular}{ccp{8cm}}
\toprule
$\nu_0$ & CV of $\Sigma_{jj}$ & Interpretation \\
\midrule
6 & $\approx 1.0$ & Minimal information; highly uncertain about $\boldsymbol{\Sigma}$. \textbf{Warning:} Severe skewness (see Section 7). \\
10 & $\approx 0.5$ & Moderate information; less skewness. \\
20 & $\approx 0.3$ & Substantial information; $\boldsymbol{\Sigma}$ fairly well-determined. \\
\bottomrule
\end{tabular}
\end{center}

\subsection{Solving for the Hyperprior Parameters}

Given the prior selection results $(\boldsymbol{\mu}_{\text{ps}}, \boldsymbol{\Sigma}_{\text{ps}})$, the variance partition ratio $\rho$, and the degrees of freedom $\nu_0$:

\textbf{Step 1: Mean hyperprior location}
\[
  \boldsymbol{\mu}_0 = \boldsymbol{\mu}_{\text{ps}}
\]

\textbf{Step 2: Shrinkage parameter (from $\rho$)}
\[
  \kappa_0 = \frac{1 - \rho}{\rho}
\]

\textbf{Step 3: Target expected population covariance}
\[
  \E[\boldsymbol{\Sigma}] = (1 - \rho) \cdot \boldsymbol{\Sigma}_{\text{ps}}
\]

\textbf{Step 4: Inverse-Wishart scale matrix}
\[
  \boldsymbol{\Psi}_0 = (\nu_0 - p - 1) \cdot \E[\boldsymbol{\Sigma}] = (\nu_0 - p - 1) \cdot (1 - \rho) \cdot \boldsymbol{\Sigma}_{\text{ps}}
\]

\subsection{Verification: The Unconditional Covariance Matches}

We verify that our construction satisfies the constraint. Starting from equation~\eqref{eq:niw_unconditional_cov}:
\begin{align*}
  \Var[\boldsymbol{\theta}_i] &= \left(1 + \frac{1}{\kappa_0}\right) \cdot \frac{\boldsymbol{\Psi}_0}{\nu_0 - p - 1} \\
  &= \left(1 + \frac{\rho}{1-\rho}\right) \cdot \frac{(\nu_0 - p - 1)(1-\rho)\boldsymbol{\Sigma}_{\text{ps}}}{\nu_0 - p - 1} \\
  &= \frac{1}{1-\rho} \cdot (1-\rho) \cdot \boldsymbol{\Sigma}_{\text{ps}} \\
  &= \boldsymbol{\Sigma}_{\text{ps}} \quad \checkmark
\end{align*}

\textbf{This identity holds for any $\rho \in (0, 1)$ and any $\nu_0 > p + 1$.} The choice of $\rho$ and $\nu_0$ affects only how the variance is decomposed and how uncertain we are about $\boldsymbol{\Sigma}$---not the unconditional distribution of individual parameters.

\subsection{What the Correlated Structure Preserves}

Unlike the independent structure, the NIW approach preserves the \textit{entire} correlation structure from prior selection:
\begin{itemize}
  \item $\E[\boldsymbol{\Sigma}] = (1-\rho) \cdot \boldsymbol{\Sigma}_{\text{ps}}$ has the same correlation structure as $\boldsymbol{\Sigma}_{\text{ps}}$
  \item $\Var[\boldsymbol{\mu}] = \frac{1}{\kappa_0} \E[\boldsymbol{\Sigma}]$ also preserves correlations
  \item The strong $a$-$c$ correlation is maintained at both population and individual levels
\end{itemize}

\subsection{Summary of the NIW Hyperprior Derivation}

\begin{center}
\begin{tabular}{ll}
\toprule
\textbf{Given} & \textbf{From} \\
\midrule
$\boldsymbol{\mu}_{\text{ps}}$, $\boldsymbol{\Sigma}_{\text{ps}}$ & Prior selection exercise (full covariance) \\
$\rho \in (0, 1)$ & Researcher choice (we use $\rho = 0.3$) \\
$\nu_0 > p + 1$ & Researcher choice (we use $\nu_0 = 6$) \\
\midrule
\textbf{Derived} & \textbf{Formula} \\
\midrule
$\boldsymbol{\mu}_0$ & $\boldsymbol{\mu}_{\text{ps}}$ \\
$\kappa_0$ & $(1 - \rho) / \rho$ \\
$\boldsymbol{\Psi}_0$ & $(\nu_0 - p - 1) \cdot (1 - \rho) \cdot \boldsymbol{\Sigma}_{\text{ps}}$ \\
\bottomrule
\end{tabular}
\end{center}

\textbf{Key property:} For any valid choice of $\rho$ and $\nu_0$, the unconditional distribution of $\boldsymbol{\theta}$ matches the prior selection target by construction.

\textbf{Comparison with independent case:} The NIW structure has only 2 researcher degrees of freedom (versus 8 for independent, or 2 with common values). This is because the matrix covariance constraint absorbs 10 degrees of freedom at once. However, the NIW structure brings computational challenges and the Inverse-Wishart skewness issue (Section 7).

\section{The Researcher's Decision}

\subsection{Constraint: Consistent Prior for a Random Individual}

We impose the constraint that the \textbf{prior predictive for a single random individual must be identical} whether we use an individual-level model or a BHM.

The prior selection exercise directly answers: ``What do credible preferences look like for a random individual?''
\[
  p(\theta) = \text{MVN}_{\text{logit}}(\boldsymbol{\mu}, \boldsymbol{\Sigma})
\]

For the BHM, the \textit{unconditional} prior (integrating over population parameters) must yield the same answer:
\[
  p_{\text{BHM}}(\theta) = \int p(\theta | \mu, \sigma^2) p(\mu) p(\sigma^2) \, d\mu \, d\sigma^2 = p(\theta)
\]

This constraint is conceptually important: both formulations encode the same prior beliefs about what a random individual's preferences look like. The difference is that the BHM explicitly decomposes this into population-level structure, allowing us to estimate \textit{how heterogeneous} the population is and to share information across individuals during estimation.

\subsection{From 16 Parameters to 2}

Whether using independent or correlated hyperpriors, the variance partition can be expressed as:
\[
  \Var[\boldsymbol{\theta}] = \underbrace{\E[\boldsymbol{\Sigma}]}_{\text{population heterogeneity}} + \underbrace{\Var[\boldsymbol{\mu}]}_{\text{uncertainty about population mean}}
\]

Two decisions determine the hyperprior structure:

\begin{enumerate}
  \item \textbf{Variance Partition Ratio} $\rho \in (0, 1)$: How to split the prior selection variance between uncertainty about the population mean and population heterogeneity.

  \item \textbf{Concentration Parameter} ($\alpha_0$ or $\nu_0$): How confident we are that the prior selection (co)variance reflects the true population (co)variance.
\end{enumerate}

These two decisions apply to both structures:

\begin{center}
\begin{tabular}{lll}
\toprule
& \textbf{Independent} & \textbf{Correlated (NIW)} \\
\midrule
Population mean prior & $\mu_j \sim N(\mu_{0,j}, \tau_{0,j}^2)$ & $\boldsymbol{\mu} \sim N(\boldsymbol{\mu}_0, \boldsymbol{\Sigma}/\kappa_0)$ \\
Population variance prior & $\sigma_j^2 \sim \text{IG}(\alpha_0, \beta_{0,j})$ & $\boldsymbol{\Sigma} \sim \text{IW}(\nu_0, \boldsymbol{\Psi}_0)$ \\
\midrule
Mean location & $\mu_{0,j} = (\boldsymbol{\mu}_{\text{logit}})_j$ & $\boldsymbol{\mu}_0 = \boldsymbol{\mu}_{\text{logit}}$ \\
Partition via $\rho$ & $\tau_{0,j}^2 = \rho \cdot \Sigma_{jj}$ & $\kappa_0 = (1-\rho)/\rho$ \\
Concentration & $\alpha_0$ (shape parameter) & $\nu_0$ (degrees of freedom) \\
Scale & $\beta_{0,j} = (1-\rho)\Sigma_{jj}(\alpha_0-1)$ & $\boldsymbol{\Psi}_0 = (1-\rho)(\nu_0-p-1)\boldsymbol{\Sigma}_{\text{logit}}$ \\
\bottomrule
\end{tabular}
\end{center}

The prior selection determines the mean and total (co)variance. The two researcher decisions ($\rho$ and concentration) then fully determine all hyperprior parameters. The correlated structure additionally preserves the full correlation matrix from prior selection.

\subsection{Decision 1: Variance Partition Ratio}

The partition ratio $\rho \in (0, 1)$ is the fraction of total prior variance attributed to uncertainty about the population mean:

\begin{center}
\begin{tabular}{cp{10cm}}
\toprule
$\rho$ & Interpretation \\
\midrule
$\rho \to 0$ & The pooled prior is wide because we expect \textit{high population heterogeneity}. The population mean is well-determined, but individuals vary greatly around it. \\
$\rho \to 1$ & The pooled prior is wide because we are \textit{uncertain about the population mean}. Once we learn the population mean, individuals cluster tightly around it. \\
$\rho = 0.3$ & (Our choice) Most prior width reflects heterogeneity (70\%), with moderate uncertainty about population location (30\%). \\
\bottomrule
\end{tabular}
\end{center}

\subsection{Decision 2: Concentration Parameter}

The concentration parameter controls how tightly the population (co)variance is distributed around its expected value:

\begin{center}
\begin{tabular}{lcc}
\toprule
& \textbf{Independent} ($\alpha_0$) & \textbf{Correlated} ($\nu_0$) \\
\midrule
Minimum for finite mean & $\alpha_0 > 1$ & $\nu_0 > p + 1 = 5$ \\
Minimum for finite variance & $\alpha_0 > 2$ & $\nu_0 > p + 3 = 7$ \\
\midrule
Low value (e.g., 6) & Diffuse; wide range of $\sigma^2$ values & Very diffuse; severe skewness \\
High value (e.g., 20) & Concentrated near $\E[\sigma^2]$ & Concentrated near $\E[\boldsymbol{\Sigma}]$ \\
\bottomrule
\end{tabular}
\end{center}

\textbf{Key difference}: The Inverse-Wishart with low $\nu_0$ exhibits extreme skewness (see Section~\ref{sec:iw-skewness}), while the Inverse-Gamma is more moderately skewed. This makes the choice of $\nu_0$ more consequential than the choice of $\alpha_0$.

Interpretation of the concentration parameter:
\begin{itemize}
  \item \textbf{Low values}: ``I am uncertain whether the prior selection variance is correct. The true population variance could be much smaller or larger.''
  \item \textbf{High values}: ``I am confident that the prior selection variance is close to the true population variance.''
\end{itemize}

\subsection{Summary: What the Researcher Must Decide}

\begin{enumerate}
  \item \textbf{Prior Selection}: Answer the question ``What do preferences look like for a single random individual?'' by defining credible RRA patterns and fitting the MVN prior. This determines $p(\theta)$ for any individual-level analysis.

  \item \textbf{Variance Partition Ratio} $\rho$: How much of the prior width comes from uncertainty about the population center vs.\ heterogeneity within the population.
  \begin{itemize}
    \item Our default: $\rho = 0.3$ (30\% location uncertainty, 70\% heterogeneity)
  \end{itemize}

  \item \textbf{Concentration Parameter} ($\alpha_0$ or $\nu_0$): How confident we are in the prior selection variance estimate.
  \begin{itemize}
    \item Independent: $\alpha_0 = 6$ gives moderate uncertainty (CV $\approx 0.5$ for $\sigma^2$)
    \item Correlated: $\nu_0 = 6$ gives minimal information but causes severe skewness; consider $\nu_0 \geq 10$ for more typical behavior
  \end{itemize}
\end{enumerate}

All other hyperprior parameters (16 total) are then \textit{derived} from the prior selection and these two decisions. The BHM decomposes the individual-level prior into population structure, but the unconditional prior for a random individual remains consistent with the prior selection.

\section{The Inverse-Wishart Skewness Problem}
\label{sec:iw-skewness}

\subsection{The Issue with Low Degrees of Freedom}

When using the Normal-Inverse-Wishart (NIW) structure for correlated hyperpriors, a critical issue arises from the choice of degrees of freedom $\nu_0$ in the Inverse-Wishart distribution.

For a $p \times p$ covariance matrix, the Inverse-Wishart distribution $\text{IW}(\nu_0, \boldsymbol{\Psi}_0)$ requires $\nu_0 > p - 1$ to be proper, and $\nu_0 > p + 1$ for the mean to exist. The mean is:
\[
  \E[\boldsymbol{\Sigma}] = \frac{\boldsymbol{\Psi}_0}{\nu_0 - p - 1}
\]

With $p = 4$ shape parameters, we need $\nu_0 > 5$ for finite mean. A natural choice is $\nu_0 = 6$, giving a denominator of $\nu_0 - p - 1 = 1$.

\subsection{Extreme Skewness of the Inverse-Wishart}

The problem is that with $\nu_0 = 6$ (just above the minimum), the Inverse-Wishart distribution is \textbf{extremely right-skewed}. While the \textit{expected value} of $\boldsymbol{\Sigma}$ matches our target, the \textit{distribution} of $\boldsymbol{\Sigma}$ values is highly asymmetric:

\begin{center}
\begin{tabular}{lc}
\toprule
Statistic & Value (trace of $\boldsymbol{\Sigma}$) \\
\midrule
Expected value & 2.125 \\
Median & 1.099 \\
5th percentile & 0.450 \\
25th percentile & 0.727 \\
75th percentile & 1.891 \\
95th percentile & 5.582 \\
\bottomrule
\end{tabular}
\end{center}

Key observations:
\begin{itemize}
  \item The \textbf{median} (1.099) is only \textbf{52\%} of the expected value (2.125)
  \item \textbf{79\%} of samples have $\text{tr}(\boldsymbol{\Sigma})$ below the expected value
  \item \textbf{48\%} of samples have $\text{tr}(\boldsymbol{\Sigma})$ less than half the expected value
  \item The mean is pulled up by rare, very large values (95th percentile is 2.6$\times$ the expected value)
\end{itemize}

\subsection{Visual Consequences}

This skewness has important consequences for the prior predictive distribution. When sampling from the NIW hierarchical model:
\begin{enumerate}
  \item Sample $\boldsymbol{\Sigma} \sim \text{IW}(\nu_0, \boldsymbol{\Psi}_0)$
  \item Sample $\boldsymbol{\mu} | \boldsymbol{\Sigma} \sim \text{MVN}(\boldsymbol{\mu}_0, \boldsymbol{\Sigma}/\kappa_0)$
  \item Sample $\boldsymbol{\theta} | \boldsymbol{\mu}, \boldsymbol{\Sigma} \sim \text{MVN}(\boldsymbol{\mu}, \boldsymbol{\Sigma})$
\end{enumerate}

Most draws of $\boldsymbol{\Sigma}$ are \textit{smaller} than expected, leading to:
\begin{itemize}
  \item Tighter clustering of $\boldsymbol{\theta}$ values around the population mean
  \item A prior predictive distribution that \textit{looks} more concentrated than the direct MVN sample from prior selection
  \item The mathematical expectation of the variance is correct, but the \textit{typical} draw has lower variance
\end{itemize}

This is not a bug---it reflects the NIW prior's ``skepticism'' about large population variance. The prior effectively says: ``I believe most populations have moderate heterogeneity, with occasional very heterogeneous populations.''

\subsection{Increasing Degrees of Freedom}

One could increase $\nu_0$ to concentrate the Inverse-Wishart closer to its mean. The coefficient of variation of $\Sigma_{jj}$ decreases as $\nu_0$ increases:
\[
  \text{CV}[\Sigma_{jj}] = \sqrt{\frac{2}{(\nu_0 - p - 1)^2(\nu_0 - p - 3)}}
\]

However, increasing $\nu_0$ changes the model's interpretation:
\begin{itemize}
  \item Higher $\nu_0$ means stronger prior beliefs about the population covariance structure
  \item With $\nu_0 = 6$, we have minimal prior information about $\boldsymbol{\Sigma}$
  \item With $\nu_0 = 20$, we have substantial prior confidence in $\E[\boldsymbol{\Sigma}]$
\end{itemize}

There is no ``correct'' choice---it depends on how confident we are that the prior selection covariance represents the true population covariance.

\section{Detailed Comparison: Independent vs.\ Correlated Hyperpriors}

We now provide a comprehensive comparison of the two hyperprior structures, including their mathematical properties, computational considerations, and practical implications.

\subsection{Mathematical Structure}

\subsubsection{Independent Hyperpriors}

Under the independent structure, each shape parameter has its own univariate hierarchical prior:
\begin{align}
  \text{logit}(\theta_{ij}) | \mu_j, \sigma_j^2 &\sim \text{Normal}(\mu_j, \sigma_j^2) \\
  \mu_j &\sim \text{Normal}(\mu_{0,j}, \tau_{0,j}^2) \\
  \sigma_j^2 &\sim \text{Inverse-Gamma}(\alpha_0, \beta_{0,j})
\end{align}
for $j \in \{a, b, c, d\}$.

The unconditional distribution of $\theta_j$ (integrating over hyperparameters) is:
\[
  \text{logit}(\theta_j) \sim t_{\nu}(\mu_{0,j}, s_j^2)
\]
where $\nu = 2\alpha_0$ and $s_j^2 = \tau_{0,j}^2 + \frac{\beta_{0,j}}{\alpha_0}$ for large $\alpha_0$, or approximately Normal when $\alpha_0$ is large.

\textbf{Key property}: Parameters are \textit{a priori independent}:
\[
  \text{Cov}[\text{logit}(\theta_j), \text{logit}(\theta_k)] = 0 \quad \text{for } j \neq k
\]

\subsubsection{Correlated Hyperpriors (Normal-Inverse-Wishart)}

Under the correlated structure:
\begin{align}
  \text{logit}(\boldsymbol{\theta}_i) | \boldsymbol{\mu}, \boldsymbol{\Sigma} &\sim \text{MVN}_4(\boldsymbol{\mu}, \boldsymbol{\Sigma}) \\
  \boldsymbol{\mu} | \boldsymbol{\Sigma} &\sim \text{MVN}_4(\boldsymbol{\mu}_0, \boldsymbol{\Sigma}/\kappa_0) \\
  \boldsymbol{\Sigma} &\sim \text{Inverse-Wishart}(\nu_0, \boldsymbol{\Psi}_0)
\end{align}

The unconditional distribution is a multivariate $t$-distribution:
\[
  \text{logit}(\boldsymbol{\theta}) \sim t_{\nu_0 - p + 1}\left(\boldsymbol{\mu}_0, \frac{\kappa_0 + 1}{\kappa_0(\nu_0 - p + 1)} \boldsymbol{\Psi}_0\right)
\]

\textbf{Key property}: Parameters can be \textit{a priori correlated}:
\[
  \text{Cov}[\text{logit}(\theta_j), \text{logit}(\theta_k)] = \frac{\kappa_0 + 1}{\kappa_0} \cdot \frac{(\boldsymbol{\Psi}_0)_{jk}}{\nu_0 - p - 1}
\]

\subsection{What Each Structure Preserves from Prior Selection}

\begin{center}
\begin{tabular}{lcc}
\toprule
Property & Independent & Correlated (NIW) \\
\midrule
Marginal means & \checkmark & \checkmark \\
Marginal variances & \checkmark & \checkmark \\
Off-diagonal covariances & $\times$ & \checkmark \\
Correlation structure & $\times$ & \checkmark \\
\midrule
Unconditional distribution shape & Approx.\ Normal & Multivariate $t$ \\
\bottomrule
\end{tabular}
\end{center}

\subsection{The Lost Correlations Problem}

The prior selection for the ATL dataset revealed strong correlations, particularly:
\[
  \text{Corr}[\text{logit}(a), \text{logit}(c)] = -0.75
\]

This correlation encodes an important constraint: utility functions with low initial curvature (small $a$) tend to have high upper curvature (large $c$), and vice versa. This maintains sensible utility shapes.

Under the independent structure, this correlation is \textit{ignored}. The prior allows combinations like (high $a$, high $c$) or (low $a$, low $c$) with equal probability, even though the prior selection identified these as implausible.

\subsection{Distributional Shape Differences}

\subsubsection{Independent: Product of Marginal $t$-Distributions}

The unconditional distribution under independent hyperpriors is:
\[
  p(\boldsymbol{\theta}) = \prod_{j=1}^{4} p(\theta_j)
\]
where each $p(\theta_j)$ is approximately a scaled $t$-distribution (or Normal for large $\alpha_0$).

This is \textit{not} a multivariate $t$-distribution---the product of univariate $t$s has lighter tails in the joint distribution than a true multivariate $t$.

\subsubsection{Correlated: Multivariate $t$-Distribution}

The NIW structure produces a true multivariate $t$-distribution with:
\begin{itemize}
  \item Heavier joint tails than the Normal
  \item Correlation between extreme values in different dimensions
  \item The skewness issues discussed in Section 7 when $\nu_0$ is small
\end{itemize}

\subsection{Typical vs.\ Expected Behavior}

This is the crucial practical difference:

\begin{center}
\begin{tabular}{lp{5cm}p{5cm}}
\toprule
& \textbf{Independent} & \textbf{Correlated (NIW)} \\
\midrule
Expected variance & Matches prior selection & Matches prior selection \\
Typical variance & Close to expected & Often \textit{smaller} than expected \\
Distribution shape & Symmetric around expectation & Right-skewed (for variance) \\
Prior predictive appearance & Similar spread to prior selection & Tighter than prior selection \\
\bottomrule
\end{tabular}
\end{center}

Under the independent structure with Inverse-Gamma($\alpha_0 = 6$, $\cdot$):
\begin{itemize}
  \item The median of $\sigma_j^2$ equals $\frac{\beta_{0,j}}{\alpha_0 + 1} \approx 0.83 \times \E[\sigma_j^2]$
  \item The distribution is moderately right-skewed
  \item Typical draws are close to the expected value
\end{itemize}

Under the correlated structure with Inverse-Wishart($\nu_0 = 6$, $\cdot$):
\begin{itemize}
  \item The median of $\text{tr}(\boldsymbol{\Sigma})$ is only $\approx 0.52 \times \E[\text{tr}(\boldsymbol{\Sigma})]$
  \item The distribution is \textit{extremely} right-skewed
  \item Typical draws are much smaller than expected
\end{itemize}

\subsection{Computational Considerations}

\begin{center}
\begin{tabular}{lcc}
\toprule
Aspect & Independent & Correlated (NIW) \\
\midrule
Number of hyperparameters to sample & 8 (4 means + 4 variances) & 14 (4 means + 10 Cholesky) \\
Conjugate updates available & \checkmark & \checkmark \\
Blocking structure & Simple (per-parameter) & Complex (matrix elements) \\
Acceptance rate sensitivity & Low & High (especially with low $\nu_0$) \\
\bottomrule
\end{tabular}
\end{center}

The NIW structure requires sampling a $4 \times 4$ covariance matrix (10 unique elements via Cholesky decomposition), leading to higher-dimensional proposal steps and potentially lower acceptance rates.

\subsection{When to Use Each Structure}

\textbf{Use Independent Hyperpriors when:}
\begin{itemize}
  \item Prior correlations are weak ($|\rho| < 0.3$)
  \item Computational efficiency is important
  \item The likelihood is highly informative (data dominates prior)
  \item You want the prior predictive to visually match the prior selection
\end{itemize}

\textbf{Use Correlated Hyperpriors (NIW) when:}
\begin{itemize}
  \item Prior correlations are strong and encode important constraints
  \item You want to model the full population covariance structure
  \item Sample sizes are very small (prior has substantial influence)
  \item You accept that the ``typical'' prior draw will be tighter than expected
\end{itemize}

\subsection{Our Pragmatic Choice}

For the Bézier utility estimation problem, we use the \textbf{independent structure} because:

\begin{enumerate}
  \item The strong $a$-$c$ correlation ($-0.75$) primarily serves to exclude implausible utility shapes. With 120 choices per subject, the likelihood will quickly rule out these shapes regardless of the prior correlation.

  \item The independent structure produces a prior predictive that visually matches the prior selection exercise, making it easier to verify and communicate.

  \item The computational simplicity allows for better MCMC mixing with smaller blocks.

  \item The NIW skewness issue means the correlated structure doesn't deliver what it promises---the \textit{typical} prior draw doesn't reflect the \textit{expected} correlation structure.
\end{enumerate}

If preserving correlations is essential, consider:
\begin{itemize}
  \item Using a higher $\nu_0$ (e.g., 15-20) to reduce skewness, at the cost of stronger prior beliefs
  \item Using an LKJ prior on the correlation matrix (separating correlation from variance)
  \item Directly parameterizing the Cholesky factor with appropriate priors
\end{itemize}

\section{Conclusion}

This document describes how BHM hyperpriors can be derived from a prior selection exercise, dramatically reducing researcher degrees of freedom:
\begin{enumerate}
  \item Using prior selection to answer: ``What do preferences look like for a single random individual?''
  \item Deriving BHM hyperpriors such that the unconditional prior matches this individual-level prior
  \item Reducing researcher degrees of freedom from 16 hyperprior parameters to 2 interpretable choices
  \item Maintaining conjugate structure (Normal-Inverse-Gamma or Normal-Inverse-Wishart) for efficient sampling
\end{enumerate}

The key insight is that the prior selection exercise answers a question about \textit{individuals}---what credible preferences look like for a random person. This answer implicitly integrates over population heterogeneity. The BHM then decomposes this integrated prior into explicit population-level components, but the unconditional prior must remain consistent with our original intuitions about individuals.

The two modeling decisions for the BHM are:
\begin{enumerate}
  \item \textbf{Variance partition ratio} $\rho$: How to split prior variance between uncertainty about the population center and heterogeneity within the population.
  \item \textbf{Concentration parameter} ($\alpha_0$ or $\nu_0$): How confident we are in the prior selection variance estimate.
\end{enumerate}

All other hyperprior parameters follow mechanically from these choices and the prior selection results.

For the correlated (NIW) structure, the concentration parameter $\nu_0$ has substantial implications due to the extreme skewness of the Inverse-Wishart distribution at low degrees of freedom. This makes the independent structure preferable in most practical applications, despite its inability to capture prior correlations.

\end{document}
